% RESEARCH_DRAFT_SHA256: 034e0cc64336d23db6ba3a0e3f5a58fec020fc2b310d386d089cbab805a38be8
% Layout and source-reference wrappers expanded for the web renderer.
\documentclass[11pt]{article}
\usepackage[margin=25mm]{geometry}
\usepackage{fontspec,unicode-math}
\setmainfont{Latin Modern Roman}
\setmathfont{Latin Modern Math}
\usepackage{amsmath,amssymb,mathtools,xurl,hyperref}
\hypersetup{colorlinks=true,urlcolor=blue,linkcolor=blue}
\setlength{\parindent}{0pt}
\setlength{\parskip}{5pt}
\setlength{\emergencystretch}{3em}
\newcommand{\sref}[2]{\hyperlink{source:#1:#2}{[S#2]}}
\newcommand{\cataloguescope}[1]{\paragraph{Source question.}#1}
\newcommand{\floor}[1]{\left\lfloor#1\right\rfloor}
\begin{document}

\section*{A determinant-norm consequence of the Agrawal congruence and its insufficiency}
\noindent catalogue entry; catalogue entry. Research attempt dated 11 October 2026.
\paragraph{Result in scope.}This manuscript contains a source-specific attempt and complete proofs of the stated auxiliary calculations. It makes no new solution or novelty claim for the entire target.
\section{Original problem and preserved definition}
The following definition and source list are reproduced from the repository. Any earlier inspection dates and reported literature claims in this preserved text describe that source record, not additional verification by this attempt.
% BEGIN PRESERVED DEFINITION
\subsection{Definitions and mathematical statement}

\paragraph{Shared notation.}Standard number systems are denoted by $\mathbb N,\mathbb Z,\mathbb Q,\mathbb R,\mathbb C$; index ranges, computational models and specialized notation follow the cited problem source and the definitions below.


For positive integers $n,r$, $(X-1)^n\equiv X^n-1\pmod{n,X^r-1}$ is equality in the quotient ring $(\mathbb Z/n\mathbb Z)[X]/(X^r-1)$. The source asks whether this congruence and $\gcd(r,n)=1$ force $n^2\equiv1\pmod r$ when $n$ is composite.
\paragraph{Statement recorded in the supplied source.}
Question 1. If (X − 1) n ≡ Xn − 1 (mod n, X r − 1) and gcd( r, n ) = 1 does this imply that

n2 ≡ 1 (mod r)when n is composite? (AKS) \href{https://aimath.org/WWN/primesinp/primesinp.pdf}{[S2]}
\paragraph{Scope qualification.}
The primary notes include heuristic discussion and conditional constructions of counterexamples; no current resolution is inferred from those historical remarks.
\paragraph{Source status.}
The supplied catalogue labels this record partially solved. Its source-status metadata was checked against the supplied snapshot on 2026-10-11. This import does not independently establish current open status or mathematical correctness.
\subsection{Short English statement}
Does the stated polynomial congruence force a composite integer to have square one modulo the polynomial period?
\subsection{Sources}
\begin{itemize}
\item[S1] Ulam AI and the UnsolvedMath Contributors, UnsolvedMath: A Curated Collection of Open Mathematics Problems, catalogue entry. Catalogue curation and metadata: CC BY 4.0. The supplied source question is retained; definitions and the English summary are editorial additions.
\url{https://huggingface.co/datasets/ulamai/UnsolvedMath}.
\item[S2] AIM Workshop Problem Lists, original problem source linked by the catalogue
\url{https://aimath.org/WWN/primesinp/primesinp.pdf}.
\end{itemize}
% END PRESERVED DEFINITION
\subsection{Verbatim supplied catalogue question}
Question 1. If (X − 1) n ≡ Xn − 1 (mod n, X r − 1) and gcd( r, n ) = 1 does this imply that

n2 ≡ 1 (mod r)when n is composite? (AKS)

\section{Attempt through a scalar invariant}
Retain positive integers $n,r$ with $n$ composite and $\gcd(n,r)=1$,
and the full hypothesis $(X-1)^n=X^n-1$ modulo $(n,X^r-1)$.
The source records $r^n\equiv r\pmod n$ as a consequence. We prove
this including nonsquarefree $n$, then test whether it can yield the
requested congruence $n^2\equiv1\pmod r$.

\paragraph{Finite-free norm proof.}
For $r=1$ the scalar conclusion is immediate. For $r\ge2$, put
$S_r(X)=1+X+\cdots+X^{r-1}$ and
$A=(\mathbb Z/n\mathbb Z)[t]/(S_r(t))$.
This is free of rank $r-1$ with basis $1,t,\ldots,t^{r-2}$,
because division by a monic polynomial works over any coefficient ring.
For $a\in A$ let $N(a)$ be the determinant of multiplication by $a$;
determinant multiplicativity gives $N(ab)=N(a)N(b)$.
Multiplication by $t$ is the companion matrix of $S_r$, so its
characteristic polynomial is $S_r$. This follows directly by expanding
$\det(XI-M_t)$ along the companion matrix columns, or equivalently
by the defining recurrence $t^{r-1}=-(1+\cdots+t^{r-2})$.
Thus $N(t-1)=(-1)^{r-1}S_r(1)=(-1)^{r-1}r$.

Since $t^r=1$, substitution $t\mapsto t^n$ preserves $S_r(t)=0$:
its $r$ terms permute those of $S_r$ as $\gcd(n,r)=1$.
An inverse substitution uses any integer $a$ with $an\equiv1\pmod r$,
so this is an algebra automorphism. Conjugating multiplication operators
by this automorphism shows $N(t^n-1)=N(t-1)$.
The given polynomial identity descends to $A$ because $S_r$ divides
$X^r-1$. Taking determinants therefore yields
\[
[(-1)^{r-1}r]^n=(-1)^{r-1}r\pmod n.
\]
The factor $r$ is a unit modulo $n$. After division the possible sign
is $(-1)^{(r-1)(n-1)}=1$: if $n$ is odd then $n-1$ is even, while
if $n$ is even then coprimality forces $r$ odd. We have proved
\[
r^{n-1}\equiv1\pmod n.
\]
No step replaced $n$ by its radical or assumed that $A$ is a field.

\paragraph{Exact obstruction to finishing with the norm alone.}
Take $n=217=7\cdot31$ and $r=5$. Since $5^6\equiv1\pmod7$
and $5^3=125\equiv1\pmod{31}$, and both $6,3$ divide $216$,
we get $5^{216}\equiv1\pmod{217}$. For the first congruence,
$5^3\equiv-1\pmod7$, so its square is one.
Yet $217^2\equiv4\pmod5$, not one. Thus the scalar norm condition
with coprimality and compositeness does not imply the desired conclusion.
The example is not a counterexample to the original conjecture:
an exact cyclic multiplication computation gives residual coefficients
\[
(X-1)^{217}-(X^{217}-1)
 \equiv207+10X^2+159X^3+58X^4\pmod{217,X^5-1},
\]
which is nonzero. The finite computation can be checked by repeated
squaring with cyclic convolution of five integer residues.

\paragraph{Automatic periods do not test the difficult implication.}
For $r$ dividing $24$, every integer coprime to $r$ has square one
modulo $r$. Indeed an odd integer $2k+1$ has square
$1+4k(k+1)\equiv1\pmod8$, and an integer not divisible by three
has square one modulo three by its two possible nonzero residues.
The Chinese remainder theorem gives the assertion for $24$, and
reduction gives it for its divisors. Thus calculations at these periods
cannot supply evidence for the nontrivial part of the conjecture.
No classification of all other periods is needed for this observation.

\section{Source context and audit}
The original Question 1 and the report's later Lenstra--Pomerance
conditional counterexample construction were inspected. No integers
satisfying all conditions of that historical construction are asserted
here. The audit checked determinant norms over composite rings,
automorphism well-definedness, signs, and the precise failure of the
scalar-only argument. Python checked the displayed residual, not the
universal conjecture. A stronger invariant retaining more polynomial
information, or an actual full-hypothesis counterexample, is still
missing. No new solution or novelty is claimed.

\section{Process and contributions}
Mathematics and model-assisted adversarial self-audit: OpenAI gpt-6-astra, ultra reasoning, as dispatched for this research batch; exact serving revision not exposed. Codex (GPT-6; exact serving revision not exposed) materially assisted with orchestration and format review. No human mathematical authorship, mathematical edits, or independent human proof checking is asserted.
Contribution prompt A was applied to the full definition and supplied source row, then prompt B was applied adversarially to the derived results. The item-specific audit above records the scope checks and rejected extrapolations. Local assembly and structural checks do not establish mathematical correctness. The source catalogue status is preserved in TOP\_PROBLEM and is separate from this attempt status.
\paragraph{Integration review.} Codex (GPT-6; exact serving revision not exposed) reviewed the written arguments and expanded layout and source-reference wrappers. Repository metadata and the production web renderer were checked using Python 3.13 and JavaScript with MathJax 3.2.2. This model-assisted review is not independent human proof verification; no PDF compilation or formal proof checking is asserted.

\section{Completion Estimate}
\noindent\textbf{COMPLETION: 0\%}
The proved material is a rederivation of known or auxiliary facts rather than a substantiated new advance on the original target. The complete scalar norm obstruction is proved, but its failure to imply the requested residue condition is explicit; no full-hypothesis counterexample or general implication is established.
\section{References}
The complete original Sources subsection above is retained, including catalogue attribution and licence notices. The additional list records primary sources actually consulted for this attempt on 11 October 2026; uncited source links reproduced above were not all rechecked.
\begin{itemize}
\item American Institute of Mathematics, Future directions in algorithmic number theory, workshop 24--28 March 2003, report version 30 April 2003; original question, relevant lecture context and adjoining remarks inspected on 11 October 2026. \url{https://aimath.org/WWN/primesinp/primesinp.pdf}.
\end{itemize}
\section{Final status}
\textbf{Unresolved in this attempt.} The complete scalar norm obstruction is proved, but its failure to imply the requested residue condition is explicit; no full-hypothesis counterexample or general implication is established.
\end{document}
